
\documentclass[final]{elsarticle}

\usepackage[margin=2.50cm]{geometry}  % tweak margins
%\usepackage{draftcopy}                % uncomment to print a "DRAFT-COPY" in background
\usepackage{graphicx}                 % if you want to include PostScript figures
\usepackage[figuresright]{rotating}   % if you have landscape tables
\usepackage{multirow}                 % fancy tables
\usepackage[toc,page]{appendix}       % appendix
\usepackage{lineno}                   % for printing the line numbers
\usepackage{natbib}                   % for building bibliography with \cite{}
\usepackage{setspace}
\usepackage{times}
\usepackage{afterpage}
\usepackage{float}
\usepackage{subcaption}
\usepackage{longtable}
\usepackage{blkarray}                 % for correlation matrices
\usepackage{color}
\usepackage{rotate}                   % rotate text in tables
\usepackage{lscape}                   % allow page to be landscape for large tables
\usepackage{amsfonts}                 % extra symbols
\usepackage{amsmath}
%\usepackage{floatpag}                 % \thisfloatpagestyle{empty}, to suppress page numbers
\usepackage{lmodern}
\usepackage{dcolumn}% Align table columns on decimal point
\usepackage{caption}% Use captionof to be able to minipage figure's \& tables together, but still get correct captions
\usepackage{enumitem}
\usepackage[breaklinks=true,hyperfigures=true]{hyperref} % make contents and references links, but do not colour them

\setenumerate{listparindent=\parindent} % Use paragraph indentation within the enumerate
\usepackage{cleveref}
% use to fix date
% \def \today { February 1st, 2009 }
% put your own definitions here:
\newcommand{\ck}{$\checkmark$}
\newcommand{\ttbs}{\char'134}
%\newcommand{\AmS}{{\protect\the\textfont2
%  A\kern-.1667em\lower.5ex\hbox{M}\kern-.125emS}}
% add words to TeX's hyphenation exception list
%\hyphenation{author another created financial paper re-commend-ed Post-Script}

\usepackage{listings}   
\lstset{language=C++}

% Particle aliases
\newcommand{\nue}{\nu_{e}}
\newcommand{\numu}{\nu_{\mu}}
\newcommand{\nutau}{\nu_{\tau}}
\newcommand{\nueb}{\xbar{\nu}_{e}}
\newcommand{\numub}{\xbar{\nu}_{\mu}}
\newcommand{\nutaub}{\xbar{\nu}_{\tau}}
\newcommand{\nub}{\xbar{\nu}}
\newcommand{\nua}{\nu_{\alpha}}
\newcommand{\nuab}{\xbar{\nu_{\alpha}}}

\newcommand{\numuany}{\overset{(-)}{\nu_\mu}}
\newcommand{\nueany}{\overset{(-)}{\nu_e}}


% By default, LaTeX doesn't like to fill more than 0.7 of a text page with tables and graphics, 
% nor does it like too many figures per page. This behaviour can be changed:
\renewcommand\floatpagefraction{.9}
\renewcommand\topfraction{.9}
\renewcommand\bottomfraction{.9}
\renewcommand\textfraction{.1}   
\setcounter{totalnumber}{50}
\setcounter{topnumber}{50}
\setcounter{bottomnumber}{50}

%define new fonts, smaller than tiny (for showing large covariance matrices in a page)
\makeatletter
\ifcase \@ptsize \relax% 10pt
  \newcommand{\miniscule}{\@setfontsize\miniscule{4}{5}}% \tiny: 5/6
  \newcommand{\verytiny}{\@setfontsize\miniscule{3}{4}}% \tiny: 5/6
\or% 11pt
  \newcommand{\miniscule}{\@setfontsize\miniscule{5}{6}}% \tiny: 6/7
  \newcommand{\verytiny}{\@setfontsize\miniscule{4}{5}}% \tiny: 6/7
\or% 12pt
  \newcommand{\miniscule}{\@setfontsize\miniscule{5}{6}}% \tiny: 6/7
  \newcommand{\verytiny}{\@setfontsize\miniscule{4}{5}}% \tiny: 6/7
\fi
\makeatother




\linenumbers

\begin{document}



{\Large \textbf{Comments on TN266v1.0 by Megan}} \\

-- It's hard to think about fitting to reconstructed energy + lepton angle,
since reconstructed energy is calculated based on the outgoing lepton angle --
how can you be sure that constraints, errors, etc aren't being double counted?
$\rightarrow$ Lines 510-513 is enough?  How about showing those results?  Are additional
bias checks needed?

\textcolor{blue}{
There is no problem with fitting Ereco and Theta if the correlations between the two variables are correctly taken into account.
Furthermore since each systematic uncertainty is treated as a different parameter, it is totally safe. 
The detector error on the lepton angle is negligible compared to the others,
so the prescription of the T2K-SK group is that the error on the lepton angle has not to be taken into account,
no matter whether Ereco o the lepton momentum is used.
So Ereco Vs Theta detector error matrix has only Ereco errors.
For what concern the cross section parameters, the normalization parameters do not depend on 
the lepton angle, while the response functions are computed as a function of Ereco and Theta.
The uncertainty on the total number of events due to the flux + cross section + Super-K-FSI-SI systematic parameters,
i.e. the systematic uncertainties which have a dependence on Theta,
are obtained by performing 10k toy experiments estimated as RMS / Mean of the distribution:
\\ \\
2D templates Ereco Vs Theta:\\ 
RHC1Re =   0.046354159\\ 
1Re =  0.041926195\\
RHC1Rmu = 0.0329263\\
1Rmu = 0.028941284\\ \\ 
1D templates Ereco: \\
RHC1Re = 0.045866498  \\
1Re =  0.041055683\\
RHC1Rmu = 0.0333726\\
1Rmu = 0.029429303\\
\\
The fractional errors are in very good agreement.
}

-- Although it isn't shown in a note, I saw plots at an analysis meeting where
this analysis had (somewhat unexpectedly) better constraints on the data fit
compared to the straight Erec analysis... Do we fully understand that
difference (since sensitivities were close to being the same)? Do we need
additional checks to understand the difference...?

\textcolor{blue}{
The agreement between Mach3, p-Theta and VALOR Asimov sensitivities is pretty good, though the templates are different.
However we saw that for Run 1-6, a particular data set, which is affected by statistical fluctuations, the 2D templates show a higher significance.
This is also seen by the analysis performed by the p-theta group.
This is due to the antineutrino mode e-like sample that in Ereco Vs Theta shows more 
events likely to be background ($\bar{\nu}_{\mu} \rightarrow \bar{\nu}_e$) rather than Ereco only (see plot below). 
This was also seen in the Run 5-6 2015 nuebar appearance analysis when the p-theta analysis was 
more compatible with null nuebar appearance with respect to both VALOR and Mach3, which were using Ereco only.
We do not think we need additional checks since we understand the source of the difference.
}

\begin{figure}[htb]
\center
\includegraphics[width=0.45\textwidth]{Ereco_ErecoVsTheta.pdf}
\end{figure}


-- line 130 : "unique" $\rightarrow$ "single"

\textcolor{blue}{DONE.}

-- line 134 : "lepton angle bin" -- what's the SK systematic error on
determining the angle (ie energy scale equivalent)?  Is that taken into account?

\textcolor{blue}{The resolution in angle is between 2-3 degrees. The prescription of T2K-SK group is to neglect it,
as it has been done so far in all the oscillation analyses involving the lepton angle.}


-- Eq1 : what's beta?

\textcolor{blue}{DONE.}


-- lines 159-161 : sentence could use a reference or additional info 

\textcolor{blue}{DONE.}


-- line 174 : "intensity" $\rightarrow$ "exposure" ? 

\textcolor{blue}{DONE.}


-- table 1, line 178-179 : "number of events" = "number of predicted events"?
w/ what oscillation parameter assumption?

\textcolor{blue}{They are oscillated only for "Oscillated" samples, which are fully oscillated. 
These numbers are taken directly from 
http://www.t2k.org/t2ksk/documentation-page/MCWeighting/MCReweighting14a}


-- line 217 + later : for this analysis, "MEC" should be changed to "2p-2h" as
far as I understand

\textcolor{blue}{DONE.}


-- lines 293-295 should go below the numu binning description

\textcolor{blue}{DONE.}

-- lines 306-307 : "The same reconstructed neutrino energy and lepton angle
binning is used for the fit of the single ring reconstructed energy spectra."
$\rightarrow$ This sentence is very hard to parse

\textcolor{blue}{DONE.}

-- lines 311-313 : this isn't true.. rec E binning hasn't changed, but binning
has certainly changed for nue sample

\textcolor{blue}{DONE.}

-- table8 : Is some additional description needed for parameters that are not
normalizations?  What does it mean when a parameter has best fit of 1.000?

\textcolor{blue}{This values are taken directly from the BANFF fit output and we would leave their full description to TN-230.
All the parameters set to 1 should correspond to the NEUT nominal value,
except CC other shape which NEUT nominal value is 0.
For all the normalization parameters, 1 is the nominal one pre-banff value.
In the text, after line 373 of v1p0, A sentence has been added to briefly explain this.}


-- line 345 : I don't think "proxies" is the right word here.. You mean
"although it's written 'numu', it's actually 'numu+nue+nutau' " ?

\textcolor{blue}{Yes, that's what I mean. The term has been changed.}


-- Figs 1 + 3 : shouldn't you be using Run 1-4 flux tuning weights for nu-mode
beam?  Or figs 1-4 are all RHC ?  (if so, that should be written somewhere.. +
how about FHC?)

\textcolor{blue}{DONE. Neutrino mode plots are added and captions are improved.}

-- tables 34-37 : need FHC, RHC labels I think

\textcolor{blue}{DONE. Added "neutrino mode" or "antineutrino mode", following the convention used in the note.}

-- line 1009 : says table 35 twice.

\textcolor{blue}{DONE.}

-- line 407-408 : do we understand that this is the cause of the bias??

\textcolor{blue}{It is not totally obvious to me because not easy to prove and understand for a such complicated analysis. 
However after the checks that were done
by the BANFF group, e.g. the flux and detector best-fit parameters of an Asimov data set are highly correlated, which could be symptom of degeneracy,
people think that probably this is the cause. Maybe for a more detailed explanation the BANFF TN is the best place to look at.   }

-- line 426 : Is it OK that the number of events from the 2015 analysis isn't
reproduced exactly?

\textcolor{blue}{Yes, as far as we understand why, which is due to different oscillation parameter values (arbitrary)
and a better cross section model.}


-- lines 434-438 : I am having trouble understanding what was done here...

\textcolor{blue}{DONE.}


-- line 473 : Probably good to define "marginalized" here..

\textcolor{blue}{I added a reference to the section where statistical methods are described. I think it could create some confusion here, since
marginalization and profiling need a quite long explanation.}


-- line 479-80 : "of about" $\rightarrow$ "by about", error reduced by how much?

\textcolor{blue}{DONE.}

-- line 481 : Is all of the run 5+6 FHC data being used for this analysis???
Maybe good to state how much POT is FHC from Run 5+6

\textcolor{blue}{DONE.}

-- fig6-17 + elsewhere : "The continuos line shows the pre-BANFF (blue) and
post-BANFF (red) ratios" $\rightarrow$ There is no continuous line (in fig6, etc) or blue
anything in any of these figures. Also "continuos" is spelled wrong "and the
oscillation parameters shown in tab. 11." $\rightarrow$ oscillation parameter priors ?

\textcolor{blue}{The caption has been improved in every plot. }

-- line 508-509 : This choice is the most conservative case for dCP = -90deg?
If so, should state that here.

\textcolor{blue}{Exactly. My understanding is that this choice is the most conservative for $\delta_{CP} = \pm \pi /2$,
since this are the values that give opposite effect to $\nu_e \rightarrow \nu_{\mu}$ and $\bar{\nu}_e \rightarrow \bar{\nu}_{\mu}$.
I maybe would avoid to state again that this is the most conservative choice for $\delta_{CP} = \pm \pi/2$,
since it is already pointed out a few sentences before. }

-- table 21 : I don't understand the column labels -- it's RMS/mean of
\#FHC/\#RHC events ?  Can you put an equation?

\textcolor{blue}{Yes it is RMS / Mean of the toy MC distribution of the ratio between neutrino mode and antineutrino mode total number of Super-K events. 
The equation has been added to the label.}

-- line 542 : "likelihood ratio method over the extended maximum-likelihood
method" $\rightarrow$ how about adding 1 sentence on what these terms mean/ what
difference is?

\textcolor{blue}{The sentence has been modified since it was misleading. By construction, this particular likelihood ratio converges to a $\chi^2$
distribution in the large sample limit and can be used for a goodness of fit test.}

-- line 652 : "+ Gaussian prior for th13"

\textcolor{blue}{DONE.}

-- table 22 : Maybe w/ reactors th13 Gaussian prior should be included in the
top table.  Isn't there a boundary problem when you constrain dcp to +/- pi
(vs +/- 2pi)?

\textcolor{blue}{It is not added to the table because the 1D confidence interval for $\sin^2 \theta_{13}$ with 
the reactor constraint is not produced, since it would be dominated by the external constraint. 
There is not boundary problem for $\delta_{CP}$. The constraint is only in the marginalization, i.e. the 
toys MC used to marginalize $\delta_{CP}$ is are produced for values between $-\pi$ and $+\pi$.
There could be an issue if the minimization algorithm, based on a descent gradient, is used.
However, as explained in section 5.3, in order to avoid a artificial local minimum at the boundary, 
4 different fits are produced, where two of them have $\pm \pi$ as starting points.
However when the full $\Delta \chi^2$ distribution is computed, this is not an issue at all. }

-- line 710-711 : This sentence is hard to understand..

\textcolor{blue}{The sentence has been improved.}

-- line 761-766 : How about not using the global minimum for the IH numbers in
this table??  Other columns wouldn't need to change + last column would be
meaningful.

\textcolor{blue}{Since the IH case is excluded at more than 90\%CL for almost all the range of parameters 
maybe it is not meaningful to show ranges with respect the IH global minimum. 
Those are 2D $1\sigma$ ranges, e.g. $\delta_{CP}$ Vs Mass Hierarchy.}


-- line 771-773 : Why isn't $dm^2_32$ fit?  Why not include the $\theta_{13}$ best fit
point?  Also, the wording is funny, since it's not a fit; it's plots of the best fit spectra

\textcolor{blue}{
The thing is for technical reason. Due to implementation of marginalization, 
it is not possible to profile more than three oscillation parameters (mass hierarchy included) in the framework.
Then it's not possible to profile all the oscillation parameters, which however are marginalized in any fit.
However best-fit spectra could have been produced also marginalizing all the parameters except $\delta_{CP}$ and mass hierarchy
or marginalizing all the oscillation parameters.
The proposal of the OA group is to profile only $\sin^2 \theta_{23}$, $\delta_{CP}$ and mass hierarchy.   
%For this reason with the OA group it was decided to marginalize $\delta m^2_{32}$ and not profile,
%since it is not degenerate with $\delta_{CP}$.
Indeed the  oscillation parameter most denerate with $\delta_{CP}$ is $\sin^2 \theta_{23}$ which is always left free and
gives a normalization to the $e$-like samples, while $\sin^2 \theta_{13}$ is almost fixed by the reactor constraint.
$\Delta m^2_{32}$ has the only effect to move the oscillation maximum to lower or higher energy
so it does not have a strong degeneracy with $\delta_{CP}$, which main effect on the spectra is simply a normalization
of the total number of events, at least with the current statistic.
}

-- line 778 : Good to explain "8-5".

\textcolor{blue}{"8-5" are exactly the values of the variable in the equation just before.
This is a quite general formula that gives the number of degrees of freedom in a fit.
Maybe the fact that "8-5" is written just after the equation could be enough. }

-- table 24 : +/- 1sigma interval for dCP doesn't match fig 23a??? This is
the main result of the analysis, I think, so it needs to be correct +
emphasized..

\textcolor{blue}{
To which $\delta_{CP}$ interval do you refer?
It seems the intervals are correct. Maybe the blue line in the T2K+reactor plot is misleading
because it is 90\% CL and not $1\sigma$ CL ($\Delta \chi^2 = 1$), the one used to set the errors. }

-- fig 24 : Need to put correct description of plot in caption

\textcolor{blue}{DONE.}

-- fig 31 : Caption is totally wrong.  Also 13b plot legend is wrong

\textcolor{blue}{DONE.}

-- line 1011 : Should separate tab 42 + description from the rest of this
sentence

\textcolor{blue}{DONE.}

-- tab 42 is hard to understand..

\textcolor{blue}{The label has been changed adding the equation.}

-- fig 37, 38 is missing pre-fit contours

\textcolor{blue}{The caption has been changed.}

-- line 1017 : Should clarify this sentence (bias is corrected by both a
correction + a new additional covariance matrix? I understood it was just a
new matrix).

\textcolor{blue}{No, also a correction to the best-fit parameters is applied.}


\newpage






{\Large \textbf{Comments on TN266v1.2 by Megan}} \\

When you make this kind of update to your TN, please put updated text in a
different color or list the updates somewhere.  It's hard to find which parts are new..
I only checked updates in Sec. 5.2-5.3 and Sec. 6.2.  Please let me know if
there are other updates.

\textcolor{blue}{There is also Appendix E, which does not concern the analysis
but only the commont method used to compute the significance.}

-- Table 23, fig 21 : Should put the true MH assumption

\textcolor{blue}{DONE.}

-- lines 755-756 : This sentence is hard to understand

\textcolor{blue}{
The point is that for higher confidence intervals the $\Delta \chi^2$ values to test are higher and 
are placed in the tail of the distribution. For example $\Delta \chi^2$ values between 0 and 2 
are likely to be included in the integral of eq. (17) for confidence intervals above 90\% CL.
The sentence has been changed to:
"In order to test confidence intervals for higher significance the tail of the $\Delta \chi^2$ distribution is more important."
}

-- line 754 : "and make the dchi2 distribution wider, as shown in fig 22" $\rightarrow$ this
isn't shown in fig 22.

\textcolor{blue}{DONE.}

-- Figure 22 + explanation : Again, I think the explanation could use some
re-wording to make it clearer.  Does it really make sense that the fit is
almost always choosing the wrong MH?  I would think that if there's a degeneracy,
the fitter should get the right MH at least ~1/2 of the time.

\textcolor{blue}{
I'm not sure we should expect 50\% of the times NH and IH,
because $\delta_{CP}$ is a continuos distribution and a small
deviation from the true value combined with systematic parameter variations, may prefer to choose IH.
I think this is favored by the still not large sensitivity on $\delta_{CP}$.
}

-- Fig 24 : again, does this flat distribution make sense?  I'd expect our data
to be able to do a little better... A table of the expected total event rates
assuming different true values of dCP and the MH could help clarify these
results.

\textcolor{blue}{DONE.}

-- line 882 : "which becomes clear when looking to eq. (19)" $\rightarrow$ I don't think
this is true.

\textcolor{blue}{DONE.}

-- eq 20 : what are f and d?

\textcolor{blue}{DONE.}

-- line 895 : "dCP=0 is favored by the data wrt dCP=pi" + next sentence $\rightarrow$ This
is backwards.

\textcolor{blue}{DONE.}


\end{document}






